% ---- replacement for the rule part of Results (from a44, resultados/analysis/44_rule_test/) ----
\subsection{Which contact positions keep their residue across lineages}

\paragraph{Positions and counting.} To avoid choosing positions after seeing their residues,
every reference position contacted (4.0~\AA) in at least 20\% of the entries with a sugar
unit at a subsite from $-3$ to $+3$ was taken: 17 positions in GH13 and 19 in GH31, the five
catalytic residues among them. Residue identity at each position was read for every chain from
the superposition (a residue belongs to a reference position when its C$\alpha$ lies within
2.0~\AA{} of the reference C$\alpha$), not only for chains whose residue touches the sugar, and
was counted once per group of proteins sharing at least 50\% sequence identity (29 groups in
GH13, 6 in GH31). Engineered mutants were removed at the mutated position (1QI5, 5TD4, 9UNS,
9XG6 at Asp289; 7WJC, 7WJF at Asp542; 3PHA, 3BLP, 3DHP, 3OLD, 3OLE, 3OLG, 3OLI, 1UA7 and 8CQF at
the positions they alter, from the \texttt{\_struct\_ref\_seq\_dif} records). The class of a
position is the contact kind made in most groups (H, side-chain hydrogen bond; Ar, apolar
contact by an aromatic residue; Al, apolar contact by an aliphatic residue). Conservation of a
position is the modal fraction
\begin{equation}
m = \frac{\max_a n_a}{\sum_a n_a},
\label{eq:modal}
\end{equation}
where $n_a$ is the number of group votes for residue type $a$ (a group whose members differ
splits its vote). For GH13 Lys182: Lys 9.5, Arg 2.5, His 2, Gln 1 of 15 groups,
$m = 9.5/15 = 0.63$. For aromatic positions the aromatic fraction (Phe, Tyr, Trp, His) is also
given.

\paragraph{What is read off.} Table~\ref{tab:rule} lists the 36 positions; 35 have a majority class
(17 H, 10 Ar, 8 Al) and one, GH31 Thr544, has none. Of the 13 hydrogen-bonding positions that contact a ring at subsite $-1$
or $+1$, eleven carry one residue type in every group ($m = 1$): in GH13 the catalytic Asp179,
Glu204 and Asp289 and the non-catalytic Arg177 and His92, with His288 at 0.97 (one natural Gln,
6SAO); in GH31 the catalytic Asp443 and Asp542 and the non-catalytic Asp327, His600, Asp203 and
Arg526. The one hydrogen-bonding position at $-1$/$+1$ that varies is the GH13 position of
Gly183 ($m = 0.50$): it is His in the eight groups where it contacts the sugar and Gly, Tyr, Glu
or Met where it does not. The ten aromatic positions keep an aromatic ring more often than
their identity (aromatic fraction 0.41--1.00, identity 0.31--1.00). The aliphatic positions vary
(GH13 Phe180 position: Ala, Val, Thr, Leu, Phe, Tyr, Met, $m = 0.24$; GH31 Ile328, $m = 0.50$),
with two exceptions: GH31 Met444 ($m = 0.83$) and GH13 Thr207, which is Lys in both of
its only two groups. At the outer subsites ($-3$, $-2$, $+2$, $+3$),
hydrogen-bonding positions vary as much as the aliphatic ones: GH13 Lys182 0.63, Met52 0.41,
Gln294 0.80, GH31 Asn207 0.50.

\paragraph{The test.} Table~\ref{tab:tests} gives the differences in mean $m$ between classes
and their significance by permutation of the class labels over positions (10\,000
permutations; 95\% intervals from 2\,000 bootstrap resamples of the sequence groups). The
hydrogen-bonding positions are more conserved than the aliphatic ones by 0.33 (all 35
positions, $p = 0.003$) and by 0.28 without the catalytic residues ($p = 0.015$). The result
holds when the class is taken from the reference structure alone instead of from the groups
($p = 0.005$--0.016), when the class is assigned on a random half of the groups and conservation
measured on the other half (positive in all 1\,000 splits; median 0.30, 95\% range
0.17--0.41), on 98 entries of 37 groups never seen with a bound sugar ($+0.31$, $p = 0.007$),
and for hydrogen-bond cut-offs of 3.2, 3.5 and 4.0~\AA{} with or without an angle criterion
($p = 0.003$--0.040). Aromatic positions keep an aromatic ring more often than aliphatic
positions keep their residue ($+0.33$, $p < 0.001$). Hydrogen-bonding and aromatic positions do
not differ significantly ($+0.18$, $p = 0.052$; held-out groups $+0.08$, $p = 0.23$). Within the
hydrogen-bonding class, positions at $-1$/$+1$ are more conserved than positions that contact
only the outer subsites (0.96 against 0.59, 13 and 4 positions, $p = 0.0015$); at the outer
subsites, hydrogen-bonding and aliphatic positions do not differ ($-0.09$, $p = 0.67$).

\paragraph{The statement the data support.} The residues whose side chains hydrogen-bond the two
rings on either side of the cut are fixed in every lineage; aromatic residues at the same
subsites are kept as aromatic rings, less strictly as the same residue; aliphatic residues at
these subsites, and residues of every kind at the outer subsites, change between lineages.
Conservation is therefore set jointly by where a residue touches the sugar and by the kind of
contact it makes, not by the kind of contact alone.

\paragraph{Limits.} GH31 rests on six groups, three of them not $\alpha$-glucosidases (an
$\alpha$-xylosidase and an $\alpha$-1,4-glucan lyase carry the Trp539$\to$Phe, Phe575$\to$His
and Met444$\to$Trp changes and, with the lyase, Ile364$\to$Asn); within the three
$\alpha$-glucosidases only Ile328, Tyr299 and five outer positions vary. Only four
hydrogen-bonding positions lie at the outer subsites. Hydrogen bonds are judged from heavy
atoms; aromatic stacking from residue type, not from ring geometry.
